% cover_letter.tex
% Cover letter for submission to
% Stochastic Environmental Research and Risk Assessment
% Compile: pdflatex serra_cover_letter.tex

\documentclass[9pt,a4paper]{article}

\usepackage[utf8]{inputenc}
\usepackage[T1]{fontenc}
\usepackage{lmodern}
\usepackage[margin=1.6cm]{geometry}
\usepackage[skip=2pt]{parskip}
\usepackage{microtype}
\usepackage{hyperref}

\hypersetup{
    colorlinks=true,
    urlcolor=blue,
    linkcolor=blue
}

\pagestyle{empty}

\begin{document}
\small

\noindent
Dear Editor-in-Chief,

I am pleased to submit my manuscript

\begin{quote}
\textbf{``When descriptive forecast-horizon gains do not imply global evidence
of incremental meteorological predictability: a multi-site PM$_{10}$ study''}
\end{quote}

for consideration for publication in
\textit{Stochastic Environmental Research and Risk Assessment}.

The manuscript is original, has not been published previously, and is not
currently under consideration by another journal.

A recurring difficulty in environmental forecasting is that a result that
looks convincing descriptively may support a much narrower conclusion once
the source of the improvement is examined statistically. This is particularly
relevant in multi-horizon forecasting. A model may outperform persistence over
a longer sequence of lead times after new predictors are added, yet the longer
skill horizon does not by itself establish that those predictors contribute
incremental predictive information. Our manuscript examines this distinction
for hourly PM$_{10}$ forecasting and asks a deliberately focused question:
when meteorological observations extend persistence-relative forecast skill,
how much evidence supports attributing that extension to the additional
meteorological information?

We study Madrid Casa de Campo and eight monitoring stations in Ireland using
an expanding-window rolling-origin evaluation over January--July 2023, with
forecast origins spaced by 24 hours and horizons from 1 to 24 hours. Persistence
and SARIMA provide temporal reference models, while the controlled
information-set comparison uses the same direct XGBoost learner first with
lagged PM$_{10}$ and then with the addition of meteorological observations
available at the forecast origin. This paired design allows the descriptive
effect of the added information to be examined without changing the learner.
It also keeps the interpretation deliberately retrospective: the
meteorological inputs are observations available at forecast issuance, not
future numerical weather prediction forecasts.

The empirical results initially suggest a strong meteorological benefit in
Madrid. Under the primary definition of the useful forecast horizon---the
longest contiguous run of positive skill relative to persistence---the horizon
increases from 9 to 17 hours when meteorological observations are added, an
8-hour extension. The corresponding relaxed last-positive-skill horizon,
however, increases from 15 to 17 hours, a 2-hour change. The contrast becomes
still less straightforward across Ireland: the mean strict horizon changes
from 21.875 to 22.875 hours (+1.000 hour), and five of the eight stations
already reach the 24-hour evaluation boundary with the lags-only model, leaving
little observable room for the strict horizon to extend further. These results make
the descriptive story interesting, but they also make clear why the number of
additional ``useful'' hours cannot by itself answer the attribution question.

That question produces the main inferential turning point of the paper. Across
the 36 planned station--horizon comparisons, all tests are valid under the
dimensionally corrected overlap-based DM-HLN specification, but none survives
global Bonferroni correction. Limited within-station signals remain under
station-wise false-discovery-rate control. We then address the most immediate
statistical objection: although the 24-hour spacing between forecast origins
implies no mechanically required positive-lag covariance from overlapping
forecast errors for the tested horizons, residual serial dependence in the
loss differential may still be present. A fixed automatic Newey--West
sensitivity analysis allows for such positive-lag dependence. The family-wise
result remains the same---again, no comparison survives global Bonferroni
correction---although the identities of the limited local signals are not
identical across the two specifications. The complete cell-level results are
supplied as Online Resource 1, alongside the main-text DM-HLN inferential
landscape (Figure 3).

For us, this distinction between the descriptive and inferential results is
the main value of the study. The Madrid skill extension remains a genuine
empirical feature of the forecasts; the inferential analysis does not make it
disappear. Rather, it changes what can reasonably be claimed from it.
Persistence-relative horizon gains, local evidence of forecast differences,
and family-wise support for incremental predictive information are different
levels of evidence. Treating them as interchangeable can make an environmental
forecasting result appear more general than the evaluation actually supports.

We believe this question is well suited to
\textit{Stochastic Environmental Research and Risk Assessment} because it
concerns not only predictive performance but the stochastic structure and
statistical interpretation of environmental forecasts. The comparison is
conducted under temporal dependence, across multiple sites and horizons, and
against persistence---a demanding reference when environmental concentrations
are strongly autocorrelated. The study also illustrates a practical
methodological consequence of the evaluation design: when forecast horizons
are expressed in physical hours but forecast origins are more widely spaced,
the dependence parameters used for forecast-comparison inference must be
defined on the loss-differential sequence rather than transferred mechanically
from physical time.

The manuscript therefore does not argue that meteorology is generally
uninformative, nor that a particular forecasting architecture is universally
superior. Its contribution is more specific: in this multi-site PM$_{10}$
experiment, substantial descriptive horizon gains can coexist with limited
family-wise evidence for attributing those gains to the added meteorological
information, and that global conclusion persists under a fixed sensitivity to
additional serial dependence. We believe that this provides a useful framework
for interpreting horizon-dependent performance claims in environmental
forecasting more generally, while keeping the empirical conclusions bounded by
the evaluated sites, period, information sets, and 24-hour forecast range.

Code, processed research materials, and figure-generation resources associated
with the study are made available through the project repository described in
the manuscript.

Thank you for considering this manuscript for publication in
\textit{Stochastic Environmental Research and Risk Assessment}. I would be
grateful for the opportunity to have the work considered by the journal's
editors and reviewers.

Yours sincerely,

\vspace{6pt}

\noindent
Federico García Crespí\\
Universidad Miguel Hernández de Elche\\
\href{mailto:fedeg@umh.es}{fedeg@umh.es}

\end{document}
